% sec-13: Document 13, the conventional pancreatic cancer trial requirements
% manual.  Operative for master prompt clause B.
% STAGE 3.  At stage 2 the closing paragraph of section 7 sat alone on a page
% carrying three lines.  The fix is the senior author's first instrument rather
% than a skip: two sentences earlier in the section are tightened, and the
% closing paragraph is shortened by one clause, so section 7 closes on the page
% it opens.  The eligibility and monitoring tables move to a page-breaking form.
\docpart{Trial Requirements}{Conventional Pancreatic Cancer Clinical Trial
Requirements Manual}

\section{Scope: What the Site Must Have Regardless of the Model}

This document is written to pass one test. A reader who deletes every large
language model and every robot from the site still holds a complete conventional
Phase 1 pancreatic cancer requirements manual. Nothing in the six sections that
follow depends on the model, and nothing in them is relaxed because the model is
present.

That order matters. A site that satisfies its model governance obligations and
falls short on investigational product accountability has failed as a clinical
trial site, and no amount of audit trail repairs it. The model earns its place
inside the conventional requirements; the conventional requirements do not bend
to accommodate it.

\section{Protocol, IND, and IRB}

The trial is an open-label, single-arm, 3+3 dose escalation of perioperative
daraxonrasib with a staged robotic pancreaticoduodenectomy in resectable
pancreatic ductal adenocarcinoma, enrolling up to 18 treated participants and up
to 30 screened \cite{phase1protocol}. The filing is the investigational new drug
application deposited in July 2026 \cite{indapp}, written against the section
skeleton of the adapted Part 312 framework \cite{cfr312adapt,cfr312}.

Three positioning statements travel with every description of this trial and are
stated here rather than left to a reader to infer. Daraxonrasib is
investigational and is already in Phase 3 evaluation; it is not described as
first in human anywhere in this package, and the supportable novelty claim
concerns the integrated surgical and advisory workflow. The robotic
configuration is specified at the site agreement and not before. No drug supply
agreement, letter of authorization, or regulatory cross-reference is in place
with the agent's developer, and no approach has been made.

Institutional review board review is obtained before any screening activity. The
board of record is the board named in the site agreement. Continuing review,
amendment submission, and reporting of unanticipated problems follow the adapted
human subject protection framework \cite{cfr50adapt,cfr50,cfr45part46}.

\section{Eligibility, Consent, and Enrollment}

\begin{mvltable}
\begin{xltabular}{\textwidth}{>{\raggedright\arraybackslash}p{2.0cm} >{\raggedright\arraybackslash}X}
\headrow \thead{Criterion} & \thead{Statement}\\
\midrule
\endfirsthead
\headrow \thead{Criterion} & \thead{Statement}\\
\midrule
\endhead
Inclusion & Histologically or cytologically confirmed pancreatic ductal adenocarcinoma \\
Inclusion & Resectable disease by multidisciplinary review, without arterial involvement \\
Inclusion & Age 18 years or older \\
Inclusion & Eastern Cooperative Oncology Group performance status 0 or 1 \\
Inclusion & Adequate hematologic, hepatic, renal, and coagulation function by protocol thresholds \\
Inclusion & Able to give written informed consent and to attend the visit schedule \\
Exclusion & Prior systemic therapy for pancreatic ductal adenocarcinoma \\
Exclusion & Borderline resectable or locally advanced disease \\
Exclusion & Distant metastatic disease on baseline imaging \\
Exclusion & Uncontrolled intercurrent illness, including active infection requiring systemic therapy \\
Exclusion & Clinically significant cardiac disease, or a corrected QT interval above the protocol threshold \\
Exclusion & Pregnancy or unwillingness to use protocol-specified contraception \\
\end{xltabular}
\end{mvltable}
\tabcapl{tab:eligibility}{Six inclusion and six exclusion criteria. The second
exclusion is the one that distinguishes this trial from the external readout it
is compared against: RASolute 302 studied metastatic, previously treated disease
and is silent on the resectable setting \cite{rasolute302}.}

Consent follows the architecture in document 02 §4 and is not restated. The
enrollment log records screening number, consent date, eligibility disposition,
assigned dose level, and, where a participant screens out, the criterion that
was not met.

\section{Dose Escalation and Dose-Limiting Toxicity}

\begin{mvtable}
\begin{tabularx}{\textwidth}{>{\raggedright\arraybackslash}p{1.3cm} >{\raggedright\arraybackslash}p{3.0cm} >{\raggedright\arraybackslash}p{2.3cm} >{\raggedright\arraybackslash}X}
\headrow \thead{Level} & \thead{Daily dose} & \thead{Participants} & \thead{Escalation rule}\\
\midrule
$-1$ & 100 mg once daily & 3 to 6 & Entered only on de-escalation from level 1 \\
1 & 160 mg once daily & 3 to 6 & 0 of 3 with a dose-limiting toxicity escalates; 1 of 3 expands to 6 \\
2 & 220 mg once daily & 3 to 6 & As level 1; 2 or more of 6 stops escalation \\
3 & 300 mg once daily & 3 to 6 & As level 1; the highest level studied \\
\end{tabularx}
\end{mvtable}
\tabcapl{tab:escalation}{Four dose levels and the escalation rule at each. Three
escalation levels at up to six participants each is the eighteen-participant
ceiling; level $-1$ is entered only by de-escalation and does not add to it.}

The dose-limiting toxicity window runs from the first dose of daraxonrasib
through postoperative day 30. A design that leaves the window unstated is not a
design, and a perioperative trial whose window ends before the operation is not
measuring the risk it exists to measure.

\begin{mvtable}
\begin{tabularx}{\textwidth}{>{\raggedright\arraybackslash}p{4.2cm} >{\raggedright\arraybackslash}X}
\headrow \thead{Grading instrument} & \thead{What it grades}\\
\midrule
CTCAE version 5.0 \cite{ctcae5} & Every non-surgical adverse event; Grade 3 or higher within the window is presumptively dose limiting \\
Clavien-Dindo \cite{claviendindo} & Surgical complications; Grade IIIb or higher within the window is presumptively dose limiting \\
ISGPS 2016 definition \cite{isgpsfistula} & Postoperative pancreatic fistula; Grade B or C is recorded and adjudicated \\
RECIST 1.1 \cite{recist11} & Radiographic response, where a target lesion is measurable \\
\end{tabularx}
\end{mvtable}
\tabcapl{tab:grading}{Four grading instruments and what each governs. Naming all
four is what allows a surgical trial to be escalated on a rule rather than on a
discussion, because a complication has a grade before it has an opinion.}

Cohort review is convened by the data manager and biostatistician within 5
business days of the last participant in a cohort completing the window. The
review record states each participant's dose-limiting toxicity status, the
grading instrument applied, and the escalation decision.

\section{Pharmacy, Laboratory, Imaging, and Biospecimens}

\textbf{Pharmacy.} Investigational product is received, stored under monitored
conditions, dispensed against a written order, reconciled at every visit, and
returned or destroyed with a record. Accountability follows the sponsor and
investigator recordkeeping requirements \cite{cfr312}. Compounding is performed
in the buffer room and the containment area built to the air change and pressure
requirements in document 08 §3, under USP General Chapters 797 and 800
\cite{usp797800}.

\textbf{Laboratory.} Safety laboratory testing is performed by a laboratory
certified under the Clinical Laboratory Improvement Amendments \cite{clia88}.
The site's own specimen processing laboratory performs collection, processing,
aliquoting, and storage only, and reports no clinical result.

\textbf{Imaging.} Baseline cross-sectional imaging is obtained within 28 days of
the first dose. Response assessment follows RECIST 1.1 \cite{recist11}.
Post-treatment imaging follows the protocol schedule. Imaging is performed off
site under the imaging agreement, which is why document 08 §4 provides no
imaging suite and no shielding.

\textbf{Biospecimens.} Tumor tissue, plasma for circulating tumor
deoxyribonucleic acid, and whole blood are collected at protocol-specified time
points. Chain of custody is recorded from collection through processing,
storage, and shipment. Specimen freezers are on the generator and carbon dioxide
backup described in document 11 §4.

\section{Monitoring, Source Documentation, Safety Reporting, and Archive}

\begin{mvtable}
\begin{tabularx}{\textwidth}{>{\raggedright\arraybackslash}p{4.0cm} >{\raggedright\arraybackslash}p{3.4cm} >{\raggedright\arraybackslash}X}
\headrow \thead{Obligation} & \thead{Interval or clock} & \thead{Authority}\\
\midrule
Site initiation visit & Before first participant & Activation gate G5, document 12 §1 \\
Monitoring visit & Every 6 weeks while enrolling; every 12 weeks otherwise & Risk-proportionate monitoring \cite{iche6r3adapt} \\
Source data verification & 100 percent for consent, eligibility, dose, and dose-limiting toxicity; risk-based otherwise & Adapted good clinical practice \cite{iche6r3adapt} \\
Serious and unexpected suspected adverse reaction & 15 calendar days & 21 CFR 312.32 \cite{cfr312} \\
Fatal or life-threatening & 7 calendar days & 21 CFR 312.32 \cite{cfr312} \\
Reportable robotic event & 5 business days; 24 hours with injury & Document 02 §6 \\
Annual report to the agency & Within 60 days of the anniversary & 21 CFR 312.33 \cite{cfr312} \\
Archive & 6 years from database lock & Document 03 §5, the longest applicable period \\
\end{tabularx}
\end{mvtable}
\tabcapl{tab:monitoring}{Eight obligations with the clock and the authority for
each. The state and federal clocks are compared in document 02 §6 and the
shorter governs; nothing in this table is the longer of two.}

Source documentation is attributable, legible, contemporaneous, original, and
accurate. A correction strikes through, initials, dates, and states a reason,
and never obscures the original entry. Electronic source records, including the
prompts and completions treated as source under document 03 §5, meet the
electronic record requirements \cite{cfr11}.

\section{The Evidence Behind the Design}

\begin{mvtable}
\begin{tabularx}{\textwidth}{>{\raggedright\arraybackslash}p{2.4cm} >{\raggedright\arraybackslash}p{1.4cm} >{\raggedright\arraybackslash}p{2.3cm} >{\raggedright\arraybackslash}p{1.4cm} >{\raggedright\arraybackslash}X}
\headrow \thead{Source} & \thead{Date} & \thead{Test arm} & \thead{Control} & \thead{Stated limitation, in the authors' own words}\\
\midrule
Empirical triplicate, 100,000 patients \cite{triplicate2025} & Jul 2025 & Grade 3+ adverse events 8.0\% & Grade 3+ 25.0\% & One log file favors experimental while the KRAS-mutant report favors control; trial-to-trial variability \\
QSP simulation, 10 arms \cite{qspmetpancre} & Aug 2025 & 12.8 months median overall survival, hazard ratio 0.25 & 5.4 months & Assumes no acquired resistance and ideal pharmacodynamics; Grade 3+ adverse events consistently high in all arms \\
Digital twin, 1000 patients \cite{daraxsims} & Oct 2025 & 12.1 months median overall survival, progression-free hazard ratio 0.31 & Not applicable & No patient-specific pharmacokinetic, pharmacodynamic, or tumor growth parameters; simple maximum-effect model, no immune compartments \\
RASolute 302 \cite{rasolute302} & May 2026 & 13.2 months median overall survival & 6.6 months & Metastatic and previously treated; silent on the resectable setting \\
\end{tabularx}
\end{mvtable}
\tabcapl{tab:evidence}{Four sources with the limitation each set of authors
stated, carried in the same row as the result. No row can be read without the
caveat that qualifies it, which is the condition on citing any of them.}

The 2025 simulated ratio of 2.4-fold and the 2026 observed ratio of 2.0-fold are
close. That proximity is a chronology observation and hypothesis-supporting. It
is \textbf{not} a validation claim, and three differences are material: 1000
simulated against 241 enrolled; a combination against a single agent; and KRAS
G12C against a primarily G12D and G12V population. Those three clauses appear
here in the same paragraph as the numbers, and they do so everywhere else in
this package.

A credibility assessment aligned to ASME V\&V 40 and to ICH M15 returned a score
of 81.9 across 55 verification notebook tests \cite{daraxsims,asmevv40,ichm15}.
That is a measured figure from the author's own verification suite and is not a
regulatory finding.

\begin{mvtable}
\begin{tabularx}{\textwidth}{>{\raggedright\arraybackslash}p{4.0cm} >{\raggedright\arraybackslash}p{3.2cm} >{\raggedright\arraybackslash}p{2.4cm} >{\raggedright\arraybackslash}X}
\headrow \thead{Work product} & \thead{Industry benchmark} & \thead{Industry time} & \thead{This author}\\
\midrule
Empirical virtual trial, triplicate & Above \$120,000 per run & 4.5 months & \$36,330 projected, about one month \\
Ten-arm QSP virtual trial & Above \$2,000,000 & Several months to over a year & \$36,330 projected, about one month \\
Digital-twin simulator platform & \$28,000 to \$700,000 & 3 to 16 weeks & \$36,330 projected, about one month \\
\end{tabularx}
\end{mvtable}
\tabcapl{tab:cost}{Three work products against their industry benchmarks. The
\$36,330 figure is a projection from the author's own source set, computed at
one operator working sixty hours a week, and is reported as projected rather
than as a completed accounting.}

A funder should read this section rather than the model governance sections,
because the decisions above are the ones a data monitoring committee will
examine: a dose-limiting toxicity window that spans the operation, four named
grading instruments, an eighteen-participant ceiling that follows from the
escalation rule rather than from a budget, and four evidence sources each
carrying its own stated limitation.
